\documentclass[10pt,a4paper]{amsart}
\usepackage[margin=19mm]{geometry}
\usepackage{amsmath,amssymb,mathtools}
\usepackage[scr=rsfs,cal=euler]{mathalfa}
\usepackage{xfrac}
\usepackage[unicode,hidelinks,bookmarks=false]{hyperref}
\newcommand{\ie}{i.e.\ }
\newcommand{\cf}{cf.\ }
\newcommand{\resp}{respectively\ }
\newcommand{\Subsection}{\subsection}
\providecommand{\nobreakdash}{\nobreak\hbox{-}}
\setlength{\emergencystretch}{2em}
\multlinegap0pt
\usepackage{bm}
\usepackage{bbm}
\usepackage{dsfont}
\usepackage{xspace}
\usepackage{cancel}
\usepackage{mathtools}
\usepackage{amsthm}
\makeatletter
\@ifundefined{theorem}{\newtheorem{theorem}{Theorem}}{}
\@ifundefined{utv}{\newtheorem{utv}{Utv}}{}
\makeatother
\title[Multidimensional quadrangle condition and cuboctahedra in la]{Multidimensional quadrangle condition and cuboctahedra in latin hypercubes}
\author{Anna A. Taranenko}
\date{Source: arXiv:2511.17082v2 (CC BY 4.0)}
\begin{document}
\maketitle

\noindent\emph{Source and licence.} Multidimensional quadrangle condition and cuboctahedra in latin hypercubes. Version arXiv:2511.17082v2 (CC BY 4.0), distributed under the Creative Commons Attribution 4.0 International licence, as recorded in the arXiv licence metadata for this exact version and shown on the abstract page https://arxiv.org/abs/2511.17082v2. Full rights notice: licenses/arxiv-2511.17082-CC-BY-4.0.txt.\par\smallskip

\noindent\textit{Editorial scope.} Counted unit: the Theorem whose source label is \texttt{lowboundsquare} in the article (source0.tex lines 234--236, source label \texttt{lowboundsquare}), delivered together with the article's own complete original proof of it (source0.tex lines 239--262, one \texttt{proof} environment of 24 lines). No mathematics is written, repaired or re-derived: statement and proof are the article's own text line for line. Required context retained uncounted: 0dimcub (lines 202--204); 1dimcub (lines 211--213); degenercub (lines 222--224). Every definition, display and label of those lines is retained unchanged. Omitted from this package and not redistributed: 1--201, 205--210, 214--221, 225--233, 237--238, 263--720. Nothing inside a retained excerpt is omitted or altered: every retained body is delivered exactly as the article's lines are written, so no omission receipt of any kind accompanies this package. Citations: the retained excerpts carry no citation command at all, so no bibliography entry of the article is transcribed and no reference is redistributed. Reference adaptation: none was needed and none was made; every reference inside the delivered unit resolves inside it, so no unresolved reference or citation is produced. Printed slips kept literal: no printed word, symbol, hypothesis or number of the counted statement or of its proof was altered. Layout and class: the article's own class is \texttt{article} with packages including amsthm, amsfonts, amsmath, amssymb, units, amsrefs, inputenc, babel, graphicx, setspace, extsizes; the wrapper is the lane's pinned \texttt{amsart} editorial wrapper at 10pt on A4, which restates the theorem-like environments the retained text needs and adds the article's own macro definitions that the retained text needs (none), with extra wrapper packages (bm, bbm, dsfont, xspace, cancel, mathtools). The delivered numbering is the unit's own. Assets: no retained span contains a figure environment, a graphics-inclusion command, a PDF-page inclusion, a table or a TikZ picture; no image file is copied into this package, the package declares no asset dependency and no image file is redistributed with it. Licence: version arXiv:2511.17082v2 of this article is distributed under the Creative Commons Attribution 4.0 International licence, as recorded in the arXiv licence metadata for this exact version and shown on the abstract page https://arxiv.org/abs/2511.17082v2. A full rights notice is delivered at licenses/arxiv-2511.17082-CC-BY-4.0.txt. Characters: every retained excerpt is pure ASCII, so the delivered engine, XeLaTeX with its default Latin Modern text font, reports no missing character.
\begin{utv} \label{0dimcub}
The number $c_0 (Q)$ of $0$-dimensional cuboctahedra in a  $d$-dimensional latin hypercube $Q$ of order $n$ is $c_0(Q) = n^{2d - 1}$.
\end{utv}

\begin{utv} \label{1dimcub}
The number $c_1 (Q)$ of $1$-dimensional cuboctahedra in a  $d$-dimensional latin hypercube $Q$ of order $n$ is $c_1(Q) =  d n^{2d-1} (n-1)$.
\end{utv}

\begin{utv} \label{degenercub}
The number $\delta_k(Q)$ of degenerate $k$-dimensional cuboctahedra in a $d$-dimensional latin hypercube $Q$ of order $n$ is $\delta_k(Q) =  {d \choose k} n^{d} (n-1)^k$.
\end{utv}

\begin{theorem} \label{lowboundsquare}
Every latin square  $L$ of order $n$ contains at least $3 n^4 + n^3 - 9 n^2 + 6n$ cuboctahedra. 
\end{theorem}

\begin{proof}
By Propositions~\ref{0dimcub} and~\ref{1dimcub}, there are $2n^4 - n^3$ cuboctahedra of dimensions $0$ and $1$  in every latin square of order $n$.  By Proposition~\ref{degenercub}, we have in addition $n^4 - 2n^3 + n^2$ degenerate cuboctahedra of dimension $2$. 

Let us estimate the number of non-degenerate $2$-dimensional cuboctahedra. 
Given symbols $x,y,z \in \{ 0, \ldots, n-1\}$,  let  $t^u_{x,y,z}$ denote the number  of $2 \times 2$ submatrices $S^u_{x,y,z}$ of the form
$$
\begin{array}{cc}
 x & y \\ 
 z & u
\end{array}
$$
 in the latin square $L$.   Since every pair of submatrices  $S^u_{x,y,z}$  in $L$ is a non-degenerate $2$-dimensional cuboctahedron, the total number of non-degenerate $2$-dimensional cuboctahedra is  equal to $\sum\limits_{x,y,z,u} t^u_{x,y,z} (t^u_{x,y,z} - 1) $. 
 From the definition of a latin square,  for all $x,y,z$, where  $x \neq y$, $x \neq z$, we have  that $\sum\limits_{u=0}^{n-1} t^u_{x,y,z} = n $ and  $ t^y_{x,y,z} =  t^z_{x,y,z} = 0$.
 
Consider the case $y \neq z$. By the pigeonhole principle, either  there exists $u$ such that $ t^u_{x,y,z}  \geq 3 $ or there exist two different $u$ and $u'$ such that $ t^u_{x,y,z}  \geq 2$ and $ t^{u'}_{x,y,z}  \geq 2$. In the first case, three submatrices $S^u_{x,y,z}$  form $6$  non-degenerate $2$-dimensional cuboctahedra, and in the second case two pairs of $S^u_{x,y,z}$ and $S^{u'}_{x,y,z}$ give  $4$  non-degenerate $2$-dimensional cuboctahedra. Assuming that for all  $x,y,z$, where $x \neq y$, $x \neq z$, $y \neq z$, the second  (worst) case is realized, we obtain at least $4 n (n-1) (n-2) = 4 n^3 - 12n^2 + 8n$  non-degenerate $2$-dimensional  cuboctahedra  of this type in total. 
 
Now assume that $y = z$. Again, by the pigeonhole principle, we see that  there exists $u$ such that $ t^u_{x,y,y}  \geq 2 $.  Since each intercalate $S^x_{x,y,y}$ is counted twice, the number $t^x_{x,y,y}$ is always even.  So in this case the minimum number of cuboctahedra is achieved when there are exactly two submatrices  $S^x_{x,y,y}$  for each $x,y$, where $x \neq y$, and all other submatrices  $S^u_{x,y,y}$ appear once. A pair of submatrices  $S^x_{x,y,y}$   gives us $2$ non-degenerate $2$-dimensional cuboctahedra, and so in total we have $2 n (n-1) = 2n^2 - 2n$ additional cuboctahedra. 
 
 Summing up, every latin square $L$ of order $n$ has at least
 $$(2n^4 - n^3) + (n^4 - 2n^3 + n^2) + (4 n^3 - 12n^2 + 8n) + (2n^2 - 2n) = 3n^4 + n^3 - 9 n^2 + 6n$$
 cuboctahedra. 
 

\end{proof}

\end{document}
